\documentclass[10pt,letterpaper]{article}

\usepackage[margin=0.45in]{geometry}
\usepackage[T1]{fontenc}
\usepackage{lmodern}
\usepackage{xcolor}
\usepackage{enumitem}
\usepackage{tabularx}
\usepackage[hidelinks]{hyperref}
\usepackage{titlesec}
\usepackage{microtype}

\definecolor{accent}{HTML}{145A78}
\setlength{\parindent}{0pt}
\setlength{\parskip}{0pt}
\setlist[itemize]{leftmargin=1.15em,itemsep=0.6pt,topsep=1pt,parsep=0pt}
\titleformat{\section}{\large\bfseries\color{accent}}{}{0pt}{}[\vspace{-2pt}\titlerule]
\titlespacing*{\section}{0pt}{4pt}{2pt}
\pagestyle{empty}

\newcommand{\role}[4]{%
  \textbf{#1}\hfill #2\\
  \textit{#3}\hfill \textit{#4}\vspace{-3pt}
}

\begin{document}
\fontsize{9}{10.4}\selectfont

\begin{center}
  {\Large\bfseries Piter Garcia}\\[1pt]
  {\large Data Scientist | Machine Learning | Applied AI}\\[2pt]
  Rochester, NY \enspace|\enspace Remote now; open to relocation after December 2026\\
  \href{mailto:pzg8794@rit.edu}{pzg8794@rit.edu}
  \enspace|\enspace +1 (646) 288-3300
  \enspace|\enspace \href{https://linkedin.com/in/garciapiterz}{linkedin.com/in/garciapiterz}\\
  \href{https://github.com/pzg8794}{github.com/pzg8794}
  \enspace|\enspace \href{https://garciapiterz.com}{garciapiterz.com}
\end{center}

\section*{Summary}
Data scientist and machine learning practitioner completing an M.S. in Data
Science in December 2026, with graduate research in AI-driven quantum network
optimization and industry experience building Python analytics, data pipelines,
and machine learning workflows in healthcare and production environments.
Experienced in model development and evaluation, statistical analysis, data
preparation, visualization, reproducible experimentation, and communicating
technical findings to varied audiences.

\section*{Technical Skills}
\begin{tabularx}{\textwidth}{@{}>{\bfseries}l X@{}}
Programming & Python, SQL, R, Java, C++, C\#, JavaScript, Bash, MATLAB\\
Data and ML & pandas, NumPy, SciPy, scikit-learn, TensorFlow, Keras, PyTorch, XGBoost, statistical modeling, feature engineering, model evaluation\\
Visualization & Matplotlib, Plotly, Seaborn, technical reporting, stakeholder presentations\\
Systems & AWS, GCP, Docker, Git/GitHub, Linux/Unix, reproducible pipelines, structured logging\\
Research & Quantum network simulation, contextual and adversarial bandits, bioinformatics, fairness-aware ML\\
\end{tabularx}

\section*{Experience}
\role{Graduate Assistant, Quantum and AI Research}{Fall 2025--Present}
{Rochester Institute of Technology, Software Engineering}{Rochester, NY}
\begin{itemize}
  \item Build a Python multi-testbed framework for threat-aware quantum routing,
    online path selection, and qubit allocation using multi-armed, contextual,
    adversarial, and neural bandit methods.
  \item Design comparison-ready experiments across stochastic and adversarial
    environments, with automated validation, structured logging, shared datasets,
    and reproducible evaluation of model utility, latency, reliability, and
    group-level outcomes.
  \item Translate large-scale experimental results into benchmark analyses,
    visualizations, technical reports, and manuscript-ready artifacts for
    faculty-led publication work.
\end{itemize}

\role{Data Solutions Engineer}{Sep 2022--Aug 2024}
{VEDADATA Inc.}{New York, NY}
\begin{itemize}
  \item Designed Python data pipelines and analytics workflows for business
    applications, emphasizing ingestion, cleaning, preprocessing, validation,
    and downstream analytics readiness.
  \item Built maintainable AWS workflows and quality checks that improved the
    traceability, reliability, and usability of production data systems.
\end{itemize}

\role{AI Data Solutions Engineer}{Jan 2021--Sep 2022}
{VIOME Inc.}{New York, NY}
\begin{itemize}
  \item Developed Python machine learning workflows for healthcare use cases,
    including preprocessing, feature engineering, supervised model
    experimentation, and evaluation.
  \item Supported reproducible model development and deployment-oriented data
    preparation in AWS-based environments.
\end{itemize}

\role{Computer Science Teacher Candidate}{2025--Present}
{University of Rochester / Pine Brook Elementary School}{Rochester, NY}
\begin{itemize}
  \item Co-design standards-aligned learning in coding, robotics, digital
    fluency, and responsible AI; use UDL supports to make technical ideas
    accessible to varied learners and collaborators.
\end{itemize}

\section*{Selected Applied Work}
\textbf{Quantum Network Path Optimization.}
Built reproducible simulation, data-processing, and evaluation workflows for
adaptive decision-making under noisy, adversarial, and resource-constrained
conditions.

\textbf{Fairness-Aware Bioinformatics and Healthcare Analytics.}
Connected hospital-readmission modeling, RNA-focused computational workflows,
and fairness assessment to evaluate predictive performance and equitable
outcomes together.

\section*{Education}
\textbf{M.S., Data Science}, Rochester Institute of Technology
\hfill Expected December 2026\\
GPA: 3.9+; focus in applied AI, bioinformatics, neural networks, and
reproducible data systems

\textbf{M.S., Teaching Computer Science K--12}, University of Rochester
\hfill Expected December 2026\\
GPA: 3.9+; NSF Robert Noyce Scholar; Advanced Certificate in Leadership in
Disability and Inclusive Practices

\textbf{M.S., Computer Science}, Rochester Institute of Technology
\hfill 2015\\
Concentrations in artificial intelligence, big data analytics, and computer
graphics

\textbf{B.S., Computer Engineering Technology}, Farmingdale State College
\hfill 2012

\section*{Recognition}
NSF Robert Noyce Scholar (2024--2026)
\enspace|\enspace IEEE Alumni Member (2009--Present)
\enspace|\enspace RIT/MESCYT merit-based graduate funding

\end{document}
